\chapter{Ice (Crystal Methamphetamine)}
\index{Ice (Crystal Methamphetamine)}

\section*{Definition and Overview}
"Ice" is the common name for crystal methamphetamine, a powerful and highly addictive stimulant drug. It appears as clear or bluish crystals that are typically smoked, but can also be injected, snorted, or swallowed. Ice produces an intense rush of euphoria, energy, and alertness by flooding the brain with dopamine and other neurotransmitters. These effects wear off quickly, and repeated use leads to a cycle of bingeing, crashing, and craving that can rapidly spiral into dependence. Ice use is associated with serious physical, psychological, and social harms, including psychosis, aggression, heart and brain damage, financial ruin, and relationship breakdown. Treatment is available and recovery is possible with the right support.

\section*{Epidemiology}
Methamphetamine is one of the most widely used illicit stimulants in many countries and the Asia-Pacific region. National surveys indicate that around 1 in 20 people have tried methamphetamine at some point, and ice is the most common form used. While the overall rate of use has remained relatively stable, patterns of use have shifted towards more frequent and dependent use among those who take the drug, and ice is a major driver of presentations to emergency departments, mental health services, and addiction treatment programs. Harms fall disproportionately on young adults, men, and communities already affected by social disadvantage, and the drug is involved in a significant proportion of violent and property crime.

\section*{Aetiology and Pathophysiology}
Methamphetamine is structurally related to the natural neurotransmitter dopamine and acts by forcing the release of large amounts of dopamine, noradrenaline, and serotonin from nerve terminals while blocking their reuptake. This produces intense stimulation of the brain's reward pathways and the sympathetic nervous system, causing the characteristic euphoria, hyperactivity, and physiological arousal. Chronic use depletes dopamine stores and damages the nerve terminals themselves, leading to tolerance, dependence, and the low, flat, and depressed mood of the withdrawal "crash". The drug also raises body temperature, heart rate, and blood pressure, and in high doses or with repeated binges can trigger psychotic symptoms, seizures, and damage to the heart, brain, and blood vessels.

\section*{Clinical Features}
\begin{itemize}
    \item Intense euphoria, energy, and confidence shortly after use, lasting several hours
    \item Reduced appetite, wakefulness, and an ability to stay awake for prolonged periods
    \item Raised heart rate, palpitations, sweating, dilated pupils, and raised blood pressure
    \item Agitation, irritability, paranoia, and suspiciousness during intoxication
    \item Teeth grinding, jaw clenching, and dry mouth
    \item The "comedown": exhaustion, depression, anxiety, and strong cravings as effects wear off
    \item With chronic use: weight loss, skin sores from picking, dental decay, and deteriorating mental health
    \item Psychosis: hallucinations, delusions, and paranoia that can persist after stopping
\end{itemize}

\section*{Differential Diagnosis}
The effects of ice can resemble other conditions, and the causes of stimulant-related harm should be considered together.
\begin{itemize}
    \item \textbf{Other stimulant drugs:} cocaine, ecstasy (MDMA), and prescription stimulants
    \item \textbf{Psychiatric disorders:} bipolar disorder, schizophrenia, and anxiety disorders presenting with agitation or psychosis
    \item \textbf{Medical emergencies:} heatstroke, stroke, seizures, and cardiac events that can mimic or complicate intoxication
    \item \textbf{Withdrawal from other substances:} alcohol, benzodiazepines, or opioids
    \item \textbf{Underlying depression or anxiety:} may be unmasked or worsened by stimulant use
\end{itemize}

\section*{When to Seek Medical Care}
Anyone who is using ice regularly, struggling to cut down, or experiencing mental or physical health problems related to the drug should seek help. GPs, drug and alcohol services, and telephone counselling lines such as the National Alcohol and Other Drug Hotline (1800 250 015) provide confidential advice and referral. People supporting a family member using ice can also access guidance and family support services.

\textbf{Contact your local emergency services for:}
\begin{itemize}
    \item A person who is unconscious, unresponsive, or not breathing
    \item Seizures, convulsions, or fitting
    \item Chest pain, palpitations, or signs of stroke, including sudden weakness or confusion
    \item Very high temperature, agitation, and confusion (possible hyperthermia or serotonin toxicity)
    \item Severe paranoia, aggression, or violence that poses a danger to the person or others
    \item Self-harm or threats of suicide
\end{itemize}

\section*{Diagnosis and Investigations}
\begin{itemize}
    \item \textbf{Clinical interview:} a non-judgemental history covering patterns of use, route, amounts, and associated harms
    \item \textbf{Physical examination:} heart rate, blood pressure, temperature, skin, teeth, and signs of injection
    \item \textbf{Urine drug screen:} can detect methamphetamine for a few days after use
    \item \textbf{Mental health assessment:} screening for depression, anxiety, psychosis, and suicide risk
    \item \textbf{Blood tests:} kidney function, muscle enzymes, and cardiac markers in severe intoxication or rhabdomyolysis
    \item \textbf{Assessment of dependence:} validated screening tools help gauge the severity of use and motivation to change
\end{itemize}

\section*{Management}
\begin{itemize}
    \item \textbf{Withdrawal management:} medically supervised support through the withdrawal phase, which typically peaks in the first week and resolves over several weeks
    \item \textbf{Cognitive behavioural therapy:} the mainstay of treatment, addressing triggers, cravings, and thinking patterns that drive use
    \item \textbf{Contingency management:} structured rewards for drug-free urine tests can help sustain abstinence
    \item \textbf{Residential rehabilitation:} structured live-in programs of several weeks to months for people with severe dependence
    \item \textbf{Mental health care:} treatment of co-occurring depression, anxiety, or psychosis, sometimes with medication
    \item \textbf{Treatment of psychosis:} antipsychotic medication and safety planning for acute or persistent stimulant psychosis
    \item \textbf{Family and peer support:} involvement of family and support networks improves outcomes
\end{itemize}

\section*{Complications}
\begin{itemize}
    \item Cardiovascular complications: raised blood pressure, arrhythmias, heart attack, and stroke, even in young people
    \item Neurological complications: seizures, cerebral haemorrhage, and permanent cognitive impairment
    \item Hyperthermia, dehydration, and rhabdomyolysis leading to kidney failure
    \item Psychiatric complications: psychosis, paranoia, aggression, depression, and suicide risk
    \item Infectious complications of injecting: HIV, hepatitis B and C, and skin and heart valve infections
    \item Dental decay, weight loss, and malnutrition
    \item Social harms: financial loss, unemployment, homelessness, relationship breakdown, and involvement with the criminal justice system
    \item Overdose and death
\end{itemize}

\section*{Prognosis and Outcomes}
Recovery from ice dependence is achievable but often takes time and multiple attempts. The withdrawal phase is unpleasant but not usually life-threatening, and the most intense cravings settle over the first few weeks. Longer-term outcomes are best with structured treatment, social support, and attention to co-occurring mental health problems. Even after stopping, some people experience persistent low mood, poor concentration, and memory difficulties that improve gradually over months. With treatment and support, many people achieve sustained abstinence or a marked reduction in use and rebuild their health, relationships, and lives.

\section*{Prevention and Self-Care}
\begin{itemize}
    \item Learn about the effects and risks of ice and talk openly with young people about drug use
    \item If using, never use alone, avoid injecting, and pace use to reduce harm
    \item Stay hydrated and take breaks from activity during any period of use
    \item Seek help early: GPs, drug and alcohol services, and helplines offer confidential, non-judgemental support
    \item Call the National Alcohol and Other Drug Hotline on 1800 250 015 for advice and referral
    \item Support friends or family members without judgment, and encourage professional help
    \item Access needle and syringe programs for sterile equipment if injecting
    \item Build healthy routines of sleep, exercise, and connection that protect against relapse
\end{itemize}

\section*{Sources and Further Reading}
The following reputable sources provide further information for healthcare students and patients seeking reliable, current advice about ice and stimulant dependence.
\begin{itemize}
    \item Alcohol and Drug Foundation. \textit{Ice and methamphetamine fact sheets.} https://adf.org.au/
    \item Healthdirect Australia. \textit{Crystal methamphetamine (ice).} https://www.healthdirect.gov.au/
    \item Cracks in the Ice. \textit{Evidence-based information about ice.} https://cracksintheice.org.au/
    \item NHS. \textit{Crystal methamphetamine: effects and getting help.} https://www.nhs.uk/
    \item MSD Manual. \textit{Amphetamine and methamphetamine intoxication.} https://www.msdmanuals.com/
    \item National Institute on Drug Abuse. \textit{Methamphetamine research report.} https://nida.nih.gov/
\end{itemize}
